% === Begin Label Data ===
% ID: 283
% 难度星级: 
% 标签: 
% 备注: 
% 组卷引用次数: 0
% === End  Label Data ===

\begin{problem}{2026}{M}{高三第一次调研测试}{5}{统计}
已知线性相关的两个变量 $x,y$ 的取值如右表所示，如果其线性回归方程为 $\hat{y}=14x-20$，那么当 $x=7$ 时的残差为
\begin{center}
\begin{tabular}{c|ccccc}
$x$ & 3 & 4 & 6 & 7 \\
\hline
$y$ & 20 & 40 & 60 & $m$
\end{tabular}
\end{center}
\begin{choices}
\choice{{2}}
\choice{{3}}
\choice{{4}}
\choice{{5}}
\end{choices}
\end{problem}

\begin{answer}
A
\end{answer}

\begin{solutions}
残差 $= y_{\text{实际}} - \hat{y}_{\text{预测}}$。  
当 $x=7$ 时，$\hat{y}=14\times7-20=98-20=78$。  
需先求 $m$。由线性回归过点 $(\bar{x},\bar{y})$：  
$\bar{x}=\dfrac{3+4+6+7}{4}=\dfrac{20}{4}=5$，  
$\bar{y}=\dfrac{20+40+60+m}{4}=\dfrac{120+m}{4}=30+\dfrac{m}{4}$。  
代入回归方程：$\bar{y}=14\bar{x}-20=14\times5-20=70-20=50$，  
故 $30+\dfrac{m}{4}=50\Rightarrow \dfrac{m}{4}=20\Rightarrow m=80$。  
则当 $x=7$ 时，实际 $y=80$，预测 $\hat{y}=78$，残差 $=80-78=2$。
\end{solutions}